\section{Antecedentes}

La obsolescencia programada es un problema que no solo afecta a los
consumidores, sino que perjudica directamente al medio ambiente debido a la
cantidad de recursos utilizados en la fabricación de productos —especialmente
electrónicos—, así como el desperdicio de muchos otros que son vistos como ``no
funcionales'' y, por lo tanto, son desechados. En 2023 se reportaron alrededor
de 6.92 mil millones de usuarios de teléfonos inteligentes en el mundo.
Considerando que el ciclo de vida promedio de estos dispositivos es de
aproximadamente 3 años, esto significa que se reemplazan constantemente, lo cual
genera una cantidad alarmante de residuos electrónicos. Además, esta cifra ni
siquiera toma en cuenta la explotación de materias primas que su producción
requiere desde el inicio de la cadena de suministro \parencite[p.~220]{ozturkcan2023}.

Según datos reportados en 2016, los Residuos de Aparatos Eléctricos y
Electrónicos (RAEE) alcanzaron aproximadamente 10.1 millones de toneladas en la
Unión Europea, lo que representó un incremento del 2.9\% respecto al año
anterior \parencite[p.~728]{malinauskaite2021}. En México, ese mismo año se
reportaron 998 mil toneladas de RAEE, siendo uno de los principales países que
generan estos residuos, ocupando el segundo lugar después de Brasil
\parencite[p.~207]{munoz2023}. Mientras la producción de equipos electrónicos
continúa en aumento, la vida útil de estos dispositivos sigue disminuyendo, lo
que inevitablemente amplifica la cantidad de residuos generados.

Se ha documentado que aproximadamente el 80\% de la contaminación y el 90\% de
los costos de fabricación asociados a los equipos electrónicos son consecuencia
de decisiones tomadas en la etapa de diseño del producto —precisamente donde
opera la obsolescencia programada \parencite[p.~728]{malinauskaite2021}. Además,
los residuos electrónicos contienen toxinas permanentes como arsénico, cadmio y
sustancias radiactivas, las cuales representan un riesgo grave para la salud de
las comunidades cercanas a los sitios de disposición, pudiendo incluso provocar
enfermedades de transmisión por el aire.

El impacto de este problema no se limita a las fronteras de los países
desarrollados. A pesar de que la exportación de residuos electrónicos con
sustancias tóxicas desde la Unión Europea hacia países no pertenecientes a la
Organización para la Cooperación y el Desarrollo Económicos (OCDE) es ilegal,
los países europeos envían cada año aproximadamente 352,474 toneladas métricas
de residuos electrónicos a naciones en vías de desarrollo, trasladando así las
consecuencias ambientales y sanitarias hacia las poblaciones más vulnerables
\parencite[p.~728]{malinauskaite2021}.

Allí, en lugares con pocos recursos y regulaciones débiles, estas comunidades
procesan los desechos de forma artesanal: incineran cables para recuperar cobre
y emplean ácidos para extraer metales preciosos de las placas. Esto contamina el
suelo, el agua y el aire, exponiendo a las personas a toxinas peligrosas. Lo que
se origina como una decisión de consumo en un país desarrollado termina
convirtiéndose en un problema de salud para comunidades vulnerables en otras
partes del mundo \parencite[p.~119]{unctad2024}.

La deshonestidad corporativa se ha intensificado con el tiempo, en parte porque
beneficia a aquellas empresas cuyo funcionamiento depende de prácticas nocivas
para la sociedad. Esto se sostiene a través de dos problemáticas: la sobrecarga
de información y el greenwashing; que permiten la persistencia de la
obsolescencia programada al aprovecharse del desinterés actual de la sociedad
para buscar información.

En estudios cualitativos realizados \parencite[p.~9]{holloway2025}, la
sobrecarga de información fue señalada por el 62\% de los participantes como uno
de los principales desafíos al momento de interpretar la extensa cantidad de
datos que las empresas divulgan sobre sus cadenas de suministro. De manera
similar, el 56\% expresó frustración ante la dificultad de acceder a información
transparente y verificable por parte de las marcas. El greenwashing tiene
relación con esta última cifra, pues ha sido señalado como una fuente principal
de escepticismo. Esta práctica —mediante la cual las empresas exageran o
tergiversan sus esfuerzos en materia de sustentabilidad— elimina la confianza
del consumidor. 

En el mercado existen aplicaciones que permiten obtener información detallada
sobre los productos:

\begin{itemize}
    \item \textbf{BARCODE LOOKUP:} Aplicación que, al escanear el código de
    barras de un producto, proporciona información técnica del mismo de forma
    inmediata, como especificaciones, características y datos relevantes del
    artículo.
    \item \textbf{GS1:} Base de datos que almacena y gestiona información
    técnica del producto junto con datos del fabricante, tales como lugar de
    fabricación, materiales utilizados en su elaboración y asociaciones con
    otras empresas.
    \item \textbf{ShopSavvy:} Aplicación móvil que permite escanear el código de
    barras de un producto, mostrando opiniones de usuarios que han adquirido
    dicho producto anteriormente.
\end{itemize}

Estas proporcionan un vistazo en el estado del mercado actual dentro del
contexto que se ha estado abordando: fuentes de información actuales que
permiten al consumidor conocer sobre los artículos que espera adquirir, así como
de sus fabricantes. 

La Gran Depresión fue un evento económico mundial originado en Estados Unidos
que provocó el inicio de varios fenómenos —como nuevas teorías económicas,
políticas macroeconómicas y cambios en instituciones económicas—. Uno de estos
fenómenos que ha influido, hasta la actualidad, en las tendencias de compra de
los consumidores es la obsolescencia programada —cuya primera aparición data en
1932 por Bernard London—como propuesta a las empresas para fomentar un mayor
consumismo por parte de la sociedad \parencite[p.~3]{corrocher2025}. La
propuesta definía un diseño deliberado de productos de tal manera que estos
tuvieran ciclos de vida más cortos con el propósito de crear una economía más
dinámica al estimular la demanda de nuevos productos, puesto que anteriormente
estos se vendían con la intención de que duraran décadas o incluso toda una vida
\parencite[p.~723]{malinauskaite2021}.

Esta práctica ha sido objeto de estudio por parte de numerosos expertos, quienes
buscan impulsar iniciativas que reduzcan el impacto ambiental y económico
generado por el consumo excesivo de productos. Entre estas iniciativas destaca el
movimiento right-to-repair —derecho a reparar—, el cual promueve que los
consumidores tengan acceso a las herramientas, piezas y conocimientos necesarios
para reparar sus propios productos, en lugar de reemplazarlos.

El movimiento right-to-repair es un movimiento moderno que ha tenido auge en
años recientes, puesto que cada vez más empresas —especialmente del sector
electrónico, industrial y automotriz— diseñan sus productos de manera que
resulten difíciles o imposibles de examinar y reparar sin herramientas
propietarias inaccesibles para el público general. Tiene como propósito
devolverle la autonomía al consumidor sobre los productos que adquiere,
exigiendo que las empresas ofrezcan acceso a piezas de repuesto a precios justos
y razonables en el mercado. En el contexto de los dispositivos electrónicos, el
movimiento también impulsa una mayor transparencia por parte de los fabricantes,
quienes deberían comunicar claramente el grado en que sus productos pueden ser
reparados. Esto permitiría a los consumidores tomar decisiones de compra más
informadas, considerando no solo el precio o las características del
dispositivo, sino también su viabilidad de reparación a largo plazo
\parencite{biscontini2024}.

Una de las primeras iniciativas dentro del movimiento right-to-repair fue el
sitio estadounidense iFixit —negocio establecido en 2003 en California,
EE.\@UU.—, el cual ofrecía manuales de reparación gratuitos especialmente para
dispositivos electrónicos. En respuesta a la creciente presión, algunas marcas
reconocidas comenzaron a establecer sus propios programas oficiales de
reparación, aunque con resultados cuestionables
\parencite[p.~220]{ozturkcan2023}.  

En 2022, Apple lanzó su programa Self-Service Repair, que otorgaba a los
usuarios acceso a documentación técnica para reparar sus dispositivos. Sin
embargo, la experiencia resultó lejos de ser conveniente para el usuario
promedio: los kits de reparación incluían piezas innecesarias para tareas
simples de reparación, así como instrucciones complejas que solo lograban
confundir a los consumidores. El resultado fue claro, parecía que la empresa no
quería que los consumidores fueran capaces de reparar por su propia cuenta,
dejando en evidencia que la iniciativa estaba lejos de representar un compromiso
real con el movimiento \parencite[p.~221]{ozturkcan2023}.

Ese mismo año, Samsung se asoció con iFixit para lanzar un programa similar que
permitía a algunos usuarios adquirir piezas de repuesto y acceder a guías de
reparación. No obstante, el programa excluía los dispositivos más recientes de
la marca y cubría únicamente un número limitado de reparaciones, lo que fue
percibido como un avance insuficiente, aunque al menos orientado en la dirección
correcta \parencite[p.~221]{ozturkcan2023}.

Estos ejemplos revelan una tendencia preocupante: mientras el discurso
corporativo adoptaba el lenguaje del movimiento, las implementaciones reales
seguían presentando barreras que, de forma sutil, continuaban incentivando al
consumidor promedio a optar por adquirir un producto nuevo antes que reparar el
que ya tenía. 

El movimiento right-to-repair es una de las iniciativas que no sólo toman en
cuenta el derecho del consumidor sobre los productos que adquiere, sino que
aborda un problema más amplio: el consumismo desmedido que caracteriza a la
sociedad actual, donde la explotación de recursos naturales, así como el
desperdicio que este provoca, es una de las muchas preocupaciones para el
desarrollo sustentable. En este sentido, el movimiento se alinea de manera
natural con los principios de la economía circular —un modelo que busca, en
palabras de \textcite[p.~229]{kirchherr2017}, alcanzar el desarrollo sustentable
mediante la creación simultánea de calidad ambiental, prosperidad económica y
equidad social, en beneficio de las generaciones presentes y futuras. La
economía circular propone cerrar el ciclo de vida de los productos mediante el
reciclaje, la reutilización y la reparación, devolviendo los bienes al mercado y
reduciendo así el impacto ambiental. 

Esta visión fue adoptada por la Comisión Europea en 2015, cuando lanzó su primer
Plan de Acción de Economía Circular bajo el principio de que ``nada debe
desperdiciarse'' buscando mejorar la economía global a través de un crecimiento
sustentable. Sin embargo, alcanzar estos objetivos resulta incompatible con la
obsolescencia programada, cuya lógica contradice directamente los fundamentos de
la economía circular, por lo que, para lograr sus objetivos principales, es
necesario primero regular —o contrarrestar— la obsolescencia programada
\parencite[p.~729]{malinauskaite2021}.

Entender el problema de la obsolescencia programada en la actualidad requiere ir
más allá de concebirla como una simple estrategia para acortar la vida útil de
los productos —el asunto es considerablemente más complejo. Los dispositivos
electrónicos no quedan obsoletos únicamente porque un componente falla. Más
bien, pierden funcionalidad por decisiones que toman los fabricantes mucho
después de que el producto fue adquirido. El problema actual no es que la
tecnología se desgaste, sino que las empresas han construido un sistema donde el
control sobre el dispositivo permanece en sus manos, incluso cuando el
consumidor ya lo pagó. 

Anteriormente, reparar un teléfono podía ser tan sencillo como cambiar una
pantalla o una batería con herramientas básicas. Sin embargo, los fabricantes
han implementado una práctica denominada “emparejamiento de piezas” o parts
pairing \parencite[p.~105]{alfieri2023}. Esto implica que cada componente
importante —como la pantalla, la batería o la cámara— posee un serial único que
debe coincidir con el de la placa principal del dispositivo. Si un usuario o
técnico independiente reemplaza una pieza, el sistema operativo la rechaza;
pudiendo mostrar mensajes de error, desactivar funciones específicas o, en casos
extremos, impedir el funcionamiento correcto del dispositivo. 

Todo este ciclo de reemplazo acelerado tiene una consecuencia muy visible: la
cantidad de residuos electrónicos crece de manera sostenida. Entre 2010 y 2022,
los desechos de pantallas, computadoras y equipos de telecomunicaciones
aumentaron un 30\% a nivel mundial, pasando de 8.1 a 10.5 millones de toneladas
\parencite[p.~108]{unctad2024}. Millones de teléfonos, computadoras son
descartados cada año, a pesar de que muchos aún podrían funcionar o ser
reparados. 

Además, la demanda de minerales necesarios para fabricar estos dispositivos
continúa en aumento. Materiales como el litio y el cobalto, esenciales para las
baterías, son cada vez más escasos y su extracción conlleva un elevado costo
ambiental y social. Se estima que para el año 2050 la necesidad de estos
minerales podría incrementarse hasta un 500\% \parencite[p.~41]{unctad2024}. En
otras palabras, se están consumiendo recursos de disponibilidad limitada para
fabricar productos que deliberadamente tienen una vida útil reducida. 

Ante este panorama, la presente investigación propuso diseñar una aplicación que
centralice datos técnicos y de sostenibilidad de productos electrónicos, con el
fin de que los usuarios puedan acceder a información real sobre la durabilidad,
costos de producción y opciones de reparación de sus dispositivos, contribuyendo
así a una toma de decisiones más informada y responsable. 

En conclusión, el consumidor requiere fuentes de información confiables que le
permitan tomar decisiones alineadas no solo con sus necesidades de consumo, sino
también con un mejor entendimiento de las prácticas corporativas que
habitualmente no son comunicadas de forma transparente. Esta información —para
que sea verdaderamente útil— debe presentarse de manera accesible, priorizando
la comprensión por encima de la exhaustividad, sin saturar al usuario con
terminología técnica o datos difíciles de interpretar. Asimismo, resulta
importante darle la oportunidad al consumidor de extender la vida útil de sus
dispositivos mediante la reutilización o la reparación, como alternativas
viables al reemplazo inmediato. 

En este contexto, si bien plataformas como BARCODE LOOKUP, GS1 o ShopSavvy son
referencias relevantes en materia de información de productos, su existencia
deja en claro un vacío tecnológico actual: la falta de una aplicación que
centralice información relevante sobre dispositivos electrónicos y sus
fabricantes, orientada específicamente a facilitar la toma de decisiones del
consumidor desde una perspectiva responsable. Es precisamente este vacío
tecnológico que el presente proyecto abordó, mediante el diseño de una
aplicación móvil que, a partir del escaneo del código de barras de un
dispositivo electrónico, permite al consumidor acceder de forma inmediata a
información sobre el producto y su fabricante, incluyendo criterios de
sostenibilidad, reputación corporativa y responsabilidad.